\documentclass[11pt,a4paper]{article}
\usepackage[utf8]{inputenc}
\usepackage[T1]{fontenc}
\usepackage{amsmath,amssymb,amsthm}
\usepackage{mathtools}
\usepackage{physics}
\usepackage{hyperref}
\usepackage{enumitem}
\usepackage{cleveref}

\title{Субквантовая теория: Часть 7 \\ Вывод Стандартной модели}
\author{На основе Теории Мод}
\date{6 июля 2026}

\newtheorem{definition}{Определение}[section]
\newtheorem{axiom}{Аксиома}[section]
\newtheorem{theorem}{Теорема}[section]
\newtheorem{lemma}{Лемма}[section]
\newtheorem{corollary}{Следствие}[section]
\newtheorem{hypothesis}{Гипотеза}[section]
\newtheorem{prediction}{Предсказание}[section]
\newtheorem{remark}{Замечание}[section]

\theoremstyle{definition}
\newtheorem{example}{Пример}[section]

\begin{document}
\maketitle
\tableofcontents
\newpage

% =================================================================
\section{От топологии мод к калибровочным группам}
% =================================================================

В Части 2 мы установили, что квантовые числа частицы — топологические инварианты её собственной моды: число нулей, их порядок, взаимное расположение. Теперь формализуем это.

\begin{definition}[Топологическое пространство моды]
Собственная мода частицы $A$ — функция $\psi_A(k) = F_{A \to A}(k)$ на дискретной оси $k = 0, 1, \dots, d_{AA}$. Рассмотрим множество нулей:
\begin{equation}
    Z_A = \{k \in [0, d_{AA}] \mid \psi_A(k) = 0\}
\end{equation}
Каждый ноль имеет кратность (порядок) $m_i$ и фазу $\phi_i$ (аргумент $\psi_A$ в окрестности нуля).
\end{definition}

\begin{definition}[Топологический заряд]
Топологический заряд нуля — целое число:
\begin{equation}
    q_i = \frac{1}{2\pi} \oint_{C_i} d(\arg \psi_A)
\end{equation}
где $C_i$ — малый контур вокруг нуля. Для изолированного нуля кратности $m$: $q_i = m \cdot \operatorname{sign}(\text{вращения фазы})$.
\end{definition}

\begin{corollary}
Сумма топологических зарядов всех нулей собственной моды равна нулю (глобальная компенсация).
\end{corollary}

Теперь свяжем топологию с калибровочными группами.

\begin{hypothesis}[Калибровочная группа как группа перестановок нулей]
Нули собственной моды с одинаковыми топологическими зарядами могут переставляться без изменения энергии сети. Группа этих перестановок — калибровочная группа частицы.
\begin{itemize}
    \item Если все нули имеют разные заряды $\to$ нет симметрии $\to$ $U(1)$ (один параметр — общая фаза).
    \item Если есть $n$ одинаковых нулей $\to$ симметрия перестановок $S_n$, которая в непрерывном пределе даёт $SU(n)$.
    \item Если есть несколько групп одинаковых нулей $\to$ прямое произведение групп.
\end{itemize}
\end{hypothesis}

% =================================================================
\section{SU(3) цвет}
% =================================================================

\begin{definition}[Цветовой заряд]
Цвет — топологический инвариант, связанный с тремя нулями собственной моды кварка, имеющими одинаковый топологический заряд $q = +1$ (или $q = -1$ для антикварков).
\end{definition}

Три одинаковых нуля $\to$ симметрия $S_3$ $\to$ $SU(3)$.

\begin{prediction}
Собственная мода кварка имеет ровно три нуля с одинаковым топологическим зарядом на основном тоне. Их взаимное расположение (расстояния между нулями) кодирует цветовое состояние.
\end{prediction}

Цветовые состояния:
\begin{itemize}
    \item Красный: нули расположены как $(0, 1, 2)$ в некоторых единицах.
    \item Зелёный: нули расположены как $(0, 1, 3)$.
    \item Синий: нули расположены как $(0, 2, 3)$.
\end{itemize}

Глюоны — моды между кварками, переставляющие эти нули. Всего $3^2 - 1 = 8$ независимых перестановок $\to$ октет глюонов.

\textbf{Конфайнмент:} В Части 2 мы вывели $C_{AB} \propto e^{-\kappa \cdot d_{AB}}$. Для кварков $\kappa_{\text{color}}$ велико $\to$ конфайнмент на масштабе $\sim 1$ ферми. Глюоны — моды между кварками — не могут существовать как свободные частицы, так как их causal history всегда привязана к кварковой паре.

% =================================================================
\section{SU(2) слабый изоспин}
% =================================================================

\begin{definition}[Слабый изоспин]
Слабый изоспин — топологический инвариант, связанный с двумя нулями собственной моды, имеющими противоположные топологические заряды: $q = +1$ и $q = -1$.
\end{definition}

Два нуля с противоположными зарядами $\to$ симметрия их совместного вращения $\to$ $SU(2)$.

\begin{prediction}
Собственная мода лептона (электрон, нейтрино) имеет два нуля с противоположными зарядами на основном тоне. Собственная мода кварка имеет три цветовых нуля плюс два слабых нуля — всего пять нулей.
\end{prediction}

Слабые бозоны ($W^+$, $W^-$, $Z$) — моды, меняющие конфигурацию этих двух нулей:
\begin{itemize}
    \item $W^+$: превращает нейтрино в электрон (сдвигает фазу одного из нулей на $\pi$).
    \item $W^-$: обратный процесс.
    \item $Z$: переставляет два нуля без изменения типа частицы.
\end{itemize}

% =================================================================
\section{U(1) гиперзаряд}
% =================================================================

\begin{definition}[Гиперзаряд]
Гиперзаряд $Y$ — общая фаза собственной моды, не связанная с нулями. Это глобальное $U(1)$-вращение:
\begin{equation}
    \psi_A \to e^{iY\theta} \psi_A
\end{equation}
\end{definition}

Электрический заряд: $Q = I_3 + Y/2$, где $I_3 = \pm 1/2$ — проекция слабого изоспина на ось $z$ (определяется относительной фазой двух слабых нулей).

Фотон — мода, соответствующая $U(1)$-вращению, не меняющая конфигурацию нулей (только общую фазу).

% =================================================================
\section{Три поколения фермионов}
% =================================================================

Вопрос: почему поколений три?

\begin{hypothesis}[Три поколения как три масштаба causal history]
Собственная мода частицы содержит нули на разных масштабах $k$. Три поколения соответствуют трём возможным конфигурациям нулей на основном тоне, первом обертоне и втором обертоне.
\begin{itemize}
    \item Основной тон ($k \approx 0 \dots K_1$): первое поколение ($e$, $\nu_e$, $u$, $d$).
    \item Первый обертон ($k \approx K_1 \dots K_2$): второе поколение ($\mu$, $\nu_\mu$, $c$, $s$).
    \item Второй обертон ($k \approx K_2 \dots K_3$): третье поколение ($\tau$, $\nu_\tau$, $t$, $b$).
\end{itemize}
\end{hypothesis}

Каждое поколение — та же топологическая конфигурация нулей, но на разных масштабах собственной моды. Масса растёт с масштабом, потому что $m_{\text{grav}} \propto \|\psi_A\|$ (норма собственной моды), а норма растёт с числом узлов.

\begin{prediction}
Существует ровно три поколения, потому что собственных обертонов, на которых топологическая конфигурация может быть стабильной, ровно три. Четвёртый обертон потребовал бы $d_{AA} >$ текущего причинного горизонта — он ещё не сформирован.
\end{prediction}

% =================================================================
\section{Калибровочные бозоны}
% =================================================================

\begin{definition}[Калибровочный бозон]
Калибровочный бозон — мода между двумя частицами, изменяющая конфигурацию нулей одной или обеих частиц на единичный шаг группы симметрии.
\end{definition}

\begin{itemize}
    \item \textbf{Фотон ($\gamma$):} $U(1)$-вращение общей фазы. Нет нулей в собственной моде $\to$ $O = 1$ (безмассовый).
    \item \textbf{$W^\pm$:} $SU(2)$-поворот двух слабых нулей. Два нуля в собственной моде $\to$ $O < 1$ (массивный).
    \item \textbf{$Z$:} $SU(2)$-поворот с изменением фазы. Два нуля $\to$ $O < 1$ (массивный).
    \item \textbf{Глюоны ($g$):} $SU(3)$-перестановка трёх цветовых нулей. Три нуля $\to$ $O < 1$ (массивные в свободном состоянии, безмассовые в конфайнменте).
\end{itemize}

Масса $W$ и $Z$ возникает из-за их собственных нулей (ортогональность $O < 1$). Механизм Хиггса в стандартной модели — эффективное описание этого факта.

% =================================================================
\section{Механизм Хиггса как рождение-исчезновение}
% =================================================================

В стандартной модели поле Хиггса даёт массу $W$ и $Z$ через спонтанное нарушение симметрии. В субквантовой теории этот механизм имеет прямую реализацию.

\begin{hypothesis}[Хиггс как частица-посредник]
Поле Хиггса — конденсат частиц-посредников $H$, постоянно рождающихся и исчезающих между $W$-бозонами и фермионами.
\end{hypothesis}

Процесс:
\begin{enumerate}
    \item $W$-бозон (мода между $A$ и $B$) имеет два нуля в собственной моде.
    \item Эти нули создают нескомпенсированность $|C_{AB}| > \theta$.
    \item Рождается частица-посредник $H$ (хиггсовский бозон) между $W$ и вакуумом.
    \item $H$ исчезает, оставляя causal history $W$ с дополнительным узлом.
    \item Этот дополнительный узел увеличивает $m_{\text{grav}}^W$ $\to$ $W$ становится массивным.
\end{enumerate}

Масса $W$:
\begin{equation}
    m_W = \Lambda \cdot (\text{число дополнительных узлов от } H\text{-рождений})^{1/2}
\end{equation}

Аналогично для $Z$. Фотон не имеет нулей $\to$ не рождает $H$ $\to$ остаётся безмассовым.

\begin{prediction}
Масса Хиггса $m_H \approx 125$ ГэВ определяется порогом $\theta$ для рождения $H$ из вакуума. Отношение $m_H / m_W \approx 125/80 \approx 1.56$ должно быть равно отношению соответствующих ортогональностей.
\end{prediction}

% =================================================================
\section{Массы фермионов и иерархия}
% =================================================================

Массы фермионов определяются их ортогональностью $O$:
\begin{equation}
    m_{\text{inert}} = \frac{S}{O}
\end{equation}
где $S$ — площадь локального участка собственной моды.

Для трёх поколений:
\begin{itemize}
    \item Первое поколение: нули на основном тоне (малые $k$) $\to$ малая $S$ $\to$ малая масса.
    \item Второе поколение: нули на первом обертоне $\to$ средняя $S$ $\to$ средняя масса.
    \item Третье поколение: нули на втором обертоне $\to$ большая $S$ $\to$ большая масса.
\end{itemize}

Иерархия масс:
\begin{equation}
    m_e : m_\mu : m_\tau \approx S_1 : S_2 : S_3 \approx 0.511 : 105.6 : 1777 \text{ МэВ}
\end{equation}

Отношения определяются отношением норм собственных мод на разных масштабах. Конкретные числа должны быть вычислены из формы поперечного спектра (провалы, плато, пики).

% =================================================================
\section{Конфайнмент и асимптотическая свобода}
% =================================================================

Из Части 2: $C_{AB} \propto e^{-\kappa \cdot d_{AB}}$.

Для сильного взаимодействия ($SU(3)$):
\begin{itemize}
    \item На малых расстояниях ($d_{AB}$ мало): корреляция велика, $\kappa_{\text{color}}$ эффективно мало $\to$ асимптотическая свобода.
    \item На больших расстояниях: корреляция экспоненциально подавлена, $\kappa_{\text{color}}$ велико $\to$ конфайнмент.
\end{itemize}

Это прямое следствие спектрального представления: для малых $d$ доминируют моды с малыми $\lambda_n$ (медленное затухание), для больших $d$ — моды с большими $\lambda_n$ (быстрое затухание).

Переход между режимами происходит при $d_{AB} \sim 1/\kappa_{\min}$, что соответствует масштабу конфайнмента $\Lambda_{\text{QCD}} \approx 200$ МэВ.

% =================================================================
\section{Электрослабое объединение}
% =================================================================

В субквантовой теории электрослабое объединение — это объединение двух типов нулей (цветовых и слабых) при высоких энергиях (малых $d$).

При $d_{AB} \to 0$:
\begin{itemize}
    \item Различие между цветовыми и слабыми нулями исчезает (все нули сливаются в один кратный ноль).
    \item $SU(3) \times SU(2) \times U(1) \to SU(5)$ или $SO(10)$ (в зависимости от точной топологии нулей).
    \item Все частицы становятся безмассовыми ($O \to 1$ при $m_{\text{grav}} \to 0$ на малых масштабах).
\end{itemize}

\begin{prediction}
При энергиях $\sim 10^{15}$ ГэВ (масштаб Великого объединения) все частицы имеют $O \approx 1$, все взаимодействия объединены. При остывании сети (расширение Вселенной) $d_{AB}$ растёт, нули разделяются, симметрия нарушается до $SU(3) \times SU(2) \times U(1)$.
\end{prediction}

% =================================================================
\section{Сводная таблица частиц Стандартной модели}
% =================================================================

\begin{center}
\begin{tabular}{c|c|c|c|c}
\textbf{Частица} & \textbf{Тип} & \textbf{Нули в собственной моде} & \textbf{O} & \textbf{Масса} \\
\hline
$\gamma$ & Бозон & $0$ нулей & $1$ & $0$ \\
$g$ & Бозон & $3$ нуля (цвет) & $<1$ & $0$ (в конфайнменте) \\
$W^\pm$ & Бозон & $2$ нуля (слабые) & $<1$ & $80.4$ ГэВ \\
$Z$ & Бозон & $2$ нуля (слабые, иная фаза) & $<1$ & $91.2$ ГэВ \\
$H$ & Бозон & $0$ нулей (вакуум) & $<1$ & $125$ ГэВ \\
$\nu_e, \nu_\mu, \nu_\tau$ & Фермион & $2$ нуля (слабые) & $\sim 1$ & $\sim 0$ \\
$e$ & Фермион & $2$ нуля (слабые) & $0.999$ & $0.511$ МэВ \\
$\mu$ & Фермион & $2$ нуля ($1$-й обертон) & $0.887$ & $105.6$ МэВ \\
$\tau$ & Фермион & $2$ нуля ($2$-й обертон) & $-0.894$ & $1777$ МэВ \\
$u, c, t$ & Кварк & $3$ нуля (цвет) $+$ $2$ нуля (слабые) & — & $\sim 2.3$ МэВ, $\sim 1.3$ ГэВ, $\sim 173$ ГэВ \\
$d, s, b$ & Кварк & $3$ нуля (цвет) $+$ $2$ нуля (слабые) & — & $\sim 4.8$ МэВ, $\sim 95$ МэВ, $\sim 4.2$ ГэВ
\end{tabular}
\end{center}

% =================================================================
\section{Открытые вопросы}
% =================================================================

\begin{enumerate}
    \item Точное вычисление масс фермионов из формы поперечного спектра (требует экспериментальных данных о шумовых паттернах).
    \item Почему слабых нулей ровно два, а цветовых — ровно три? (Гипотеза: определяется размерностью сети на момент декаплинга.)
    \item Смешивание поколений (матрица CKM) — как перестановки нулей между обертонами.
    \item CP-нарушение — как асимметрия фаз нулей на разных обертонах.
\end{enumerate}

% =================================================================
\section*{Итог Части 7}
% =================================================================

Стандартная модель возникает из субквантовой теории как прямое следствие топологии нулей собственных мод:

\begin{itemize}
    \item $SU(3)$ — перестановки трёх цветовых нулей.
    \item $SU(2)$ — вращения двух слабых нулей с противоположными зарядами.
    \item $U(1)$ — общая фаза.
    \item Три поколения — три масштаба (обертона) собственной моды.
    \item Калибровочные бозоны — моды, меняющие конфигурацию нулей.
    \item Хиггс — частица-посредник, рождающаяся и исчезающая между $W/Z$ и вакуумом.
    \item Массы — через ортогональность $O$.
    \item Конфайнмент — экспоненциальное подавление корреляции с расстоянием.
\end{itemize}

Всё без дополнительных предположений — только сеть, моды, спектры $F$, ортогональность $O$ и топология нулей.

\end{document}